\chapter{TỔNG QUAN DỰ ÁN VÀ PHÂN TÍCH YÊU CẦU}
\label{ch:chuong1}

\section{Giới thiệu và bối cảnh đề tài}
\label{sec:gioi_thieu}

Trong bối cảnh chuyển đổi số giáo dục đại học và phát triển các mô hình khuôn viên thông minh (Smart Campus), hệ thống ký túc xá (KTX) sinh viên giữ vai trò là hạ tầng sinh hoạt và học tập thiết yếu phục vụ hàng ngàn sinh viên nội trú. Một môi trường nội trú an toàn, văn minh và tiện nghi không chỉ bảo đảm sức khỏe, tính mạng cho sinh viên mà còn nâng cao hiệu quả quản lý cơ sở vật chất của nhà trường.

Tuy nhiên, qua khảo sát thực tế quy trình vận hành tại các ký túc xá sinh viên hiện nay, công tác quản trị vẫn bộc lộ nhiều điểm bất cập và rủi ro nghiêm trọng:
\begin{itemize}
    \item \textbf{Hiểm họa cháy nổ và mất an toàn phòng cháy chữa cháy (PCCC):} Nguy cơ hỏa hoạn do sự cố rò rỉ khí gas tại khu vực nấu ăn hoặc chập điện các thiết bị sinh hoạt công suất cao là mối đe dọa thường trực. Hầu hết các ký túc xá hiện nay chỉ trang bị chuông báo độc lập tại chỗ hoặc dựa vào tuần tra thủ công, thiếu cơ chế phát hiện sớm nồng độ khí gas nguy hiểm, không có khả năng tự động ngắt nguồn điện phòng để triệt tiêu nguồn phát tia lửa điện, và không thể gửi cảnh báo khẩn cấp tức thời đến ban quản lý.
    \item \textbf{Kiểm soát an ninh ra vào thủ công và thiếu chính xác:} Việc quản lý sinh viên ra vào phòng phần lớn vẫn dựa vào chìa khóa cơ truyền thống hoặc sổ sách trực ban. Chìa khóa cơ dễ bị sao chép trái phép, dễ thất lạc, gây khó khăn cho việc quản lý giờ giấc và không thể truy vết khi xảy ra sự cố mất cắp tài sản.
    \item \textbf{Thiếu vắng nền tảng số hóa quản trị tập trung:} Chưa có một hệ sinh thái phần mềm kết nối trực tiếp giữa ban quản lý và sinh viên. Khi xảy ra sự cố phòng ở, sinh viên phải báo cáo trực tiếp bằng giấy tờ hoặc liên hệ thủ công, dẫn đến độ trễ xử lý cao, thiếu dữ liệu lịch sử để thống kê và tối ưu hóa vận hành.
\end{itemize}

Nhằm giải quyết triệt để các bài toán thực tiễn nêu trên, đề tài \textbf{``Hệ thống Quản lý Ký túc xá Thông minh (SmartDorm)''} được nghiên cứu và xây dựng toàn diện, kết hợp chặt chẽ giữa phần cứng nhúng Internet of Things (IoT), hệ thống máy chủ backend xử lý thời gian thực, cơ sở dữ liệu quan hệ, phân tích trí tuệ nhân tạo (AI Anomaly Detection \& Trợ lý ảo Copilot) và giao diện web hiện đại phân quyền theo vai trò.

\section{Phân tích yêu cầu hệ thống}
\label{sec:phan_tich_yeu_cau}

\subsection{Yêu cầu chức năng}
\label{subsec:yeu_cau_chuc_nang}

Hệ thống SmartDorm phải đáp ứng đầy đủ các nhóm chức năng nghiệp vụ cốt lõi sau:
\begin{enumerate}
    \item \textbf{Giám sát môi trường và vi khí hậu phòng ở:}
    \begin{itemize}
        \item Thu thập liên tục các chỉ số nhiệt độ ($^\circ\text{C}$) và độ ẩm (\%) không khí phòng ở thông qua cảm biến kỹ thuật số DHT22 với chu kỳ từ 3 đến 5 giây.
        \item Cập nhật dữ liệu tức thời lên bảng điều khiển thời gian thực (Real-time Dashboard) mà không yêu cầu người dùng phải tải lại trang.
    \end{itemize}
    \item \textbf{Giám sát an toàn PCCC và phản ứng khẩn cấp tự động:}
    \begin{itemize}
        \item Đo lường liên tục nồng độ khí gas và khói thông qua cảm biến chuyên dụng MQ-2.
        \item Tích hợp bộ quy tắc xử lý an toàn (Rule Engine): Khi nồng độ khí gas đạt ngưỡng cảnh báo ($\ge 50\%$) hoặc nhiệt độ bất thường ($> 45^\circ\text{C}$), kích hoạt còi buzzer bíp ngắt quãng và gửi cảnh báo mức WARNING.
        \item Khi nồng độ khí gas đạt ngưỡng nguy hiểm cháy nổ ($\ge 80\%$ hoặc sự cố FIRE): Ngay lập tức kích hoạt còi hú liên tục, \textbf{tự động ngắt toàn bộ relay điện phòng} nhằm triệt tiêu tia lửa điện, đồng thời phát lệnh khẩn cấp (Emergency Broadcast) gửi thông báo đa kênh qua Telegram Bot, Email SMTP và pop-up cảnh báo đỏ trên Web.
    \end{itemize}
    \item \textbf{Kiểm soát an ninh ra vào bằng thẻ từ RFID và thuật toán tự động:}
    \begin{itemize}
        \item Quản lý lượt ra vào phòng ở thông qua công nghệ thẻ từ RFID chuẩn 13.56 MHz (RC522).
        \item Ứng dụng thuật toán sinh mã thẻ tự động dựa trên hàm băm SHA-256 kết hợp mã số sinh viên (MSSV) để cấp phát mã thẻ chuẩn định dạng \texttt{CARD\_XXXXXXXX}.
        \item Tự động đối soát UID thẻ với danh sách sinh viên nội trú của phòng: Nếu hợp lệ (\texttt{GRANTED}), kích hoạt mở khóa cửa Solenoid trong 4 giây và còi bíp xác nhận; nếu không hợp lệ (\texttt{DENIED}), giữ khóa cửa, phát 3 tiếng còi cảnh báo và ghi nhận vi phạm an ninh lên hệ thống.
    \end{itemize}
    \item \textbf{Điều khiển thiết bị phòng ở và cơ chế hoàn tác lệnh (Undo):}
    \begin{itemize}
        \item Điều khiển bật/tắt an toàn các tải thiết bị phòng: Đèn chiếu sáng, Quạt thông gió và Khóa cửa điện từ thông qua module relay cách ly quang và các nút bấm vật lý trên bo mạch.
        \item Áp dụng mẫu thiết kế Command Pattern ghi nhận lịch sử mọi thao tác điều khiển (ai thực hiện, thời gian nào, qua web hay nút bấm), cung cấp tính năng Hoàn tác lệnh gần nhất (Undo) để khôi phục trạng thái thiết bị tức thời.
    \end{itemize}
    \item \textbf{Phân tích thông minh AI và Trợ lý ảo AI Copilot:}
    \begin{itemize}
        \item Phát hiện bất thường môi trường (AI Anomaly Detection) dựa trên thuật toán Isolation Forest nhận diện các chuỗi gia tăng nhiệt độ hoặc khí gas bất thường.
        \item Tích hợp trợ lý ảo AI Copilot hỗ trợ giải thích ngữ cảnh sự cố, đề xuất phương án xử lý theo quy chuẩn an toàn và hỏi đáp dữ liệu phòng ở thời gian thực.
    \end{itemize}
    \item \textbf{Báo cáo thống kê và xuất dữ liệu thật:}
    \begin{itemize}
        \item Xuất dữ liệu báo cáo chuyên sâu định dạng Excel và CSV phục vụ công tác thanh tra và lưu trữ: Báo cáo vi khí hậu môi trường, Báo cáo kiểm soát an ninh ra vào RFID, Báo cáo sự cố khẩn cấp PCCC và Báo cáo nhật ký thiết bị.
    \end{itemize}
\end{enumerate}

\subsection{Yêu cầu an toàn và bảo mật}
\label{subsec:yeu_cau_bao_mat}

\begin{itemize}
    \item \textbf{Phân quyền truy cập dựa trên vai trò (RBAC):} Phân tách quyền hạn rõ rệt giữa hai phân hệ Quản trị viên (Admin Portal) và Sinh viên nội trú (Student Portal). Sinh viên chỉ có quyền truy cập dữ liệu và điều khiển thiết bị thuộc phòng ở của mình.
    \item \textbf{Bảo mật phiên làm việc và API:} Mọi giao tiếp giữa trình duyệt web và máy chủ backend đều được xác thực thông qua JSON Web Token (JWT) với mã hóa mật khẩu an toàn theo giải thuật băm Bcrypt Salt.
    \item \textbf{An toàn điện học đường tuyệt đối:} Phần cứng thử nghiệm tuân thủ quy chuẩn điện hạ áp, sử dụng nguồn 5V DC/3.3V DC an toàn qua cổng Type-C, loại bỏ hoàn toàn việc can thiệp trực tiếp vào điện lưới cao áp 220V nguy hiểm.
    \item \textbf{Ngăn chặn truy cập trái phép phần cứng:} Module mạng IoT ESP32 kết nối MQTT Broker qua kênh xác thực người dùng và mật khẩu riêng biệt.
\end{itemize}

\subsection{Yêu cầu phi chức năng và hiệu năng}
\label{subsec:yeu_cau_hieu_nang}

\begin{itemize}
    \item \textbf{Độ trễ truyền nhận vi điều khiển (MQTT Latency):} Thời gian từ khi cảm biến phần cứng phát tín hiệu đến khi máy chủ Backend nhận và xử lý gói tin qua MQTT Broker không vượt quá $50\text{ ms}$ (thực tế đo đạc đạt dưới $40\text{ ms}$).
    \item \textbf{Độ trễ truyền phát thời gian thực (WebSocket Latency):} Thời gian cập nhật trạng thái từ backend lên giao diện trình duyệt web qua Socket.IO đạt dưới $15\text{ ms}$.
    \item \textbf{Thời gian phản ứng ngắt điện PCCC khẩn cấp:} Từ thời điểm nồng độ khí gas vượt ngưỡng nguy hiểm $\ge 80\%$, toàn bộ quy trình gồm: hú còi tại chỗ, ngắt relay điện phòng, cập nhật trạng thái lên cơ sở dữ liệu và gửi thông báo Telegram phải hoàn tất trong vòng dưới $1.5\text{ giây}$.
    \item \textbf{Khả năng chịu tải đồng thời:} Hệ thống duy trì hoạt động ổn định khi phục vụ đồng thời 500 kết nối WebSockets và 100 yêu cầu truy vấn API mỗi giây mà không bị nghẽn đường truyền.
\end{itemize}

\section{Thiết kế hệ thống tổng thể}
\label{sec:thiet_ke_tong_the}

Hệ thống SmartDorm được tổ chức theo kiến trúc phân tầng 5 lớp (5-Tier Architecture), bảo đảm tính tách biệt độc lập giữa các thành phần, dễ dàng bảo trì và mở rộng linh hoạt:

\begin{figure}[H]
\centering
\begin{tikzpicture}[
    layer/.style={rectangle, draw=blue!80!black, fill=blue!5, thick, minimum width=14.5cm, minimum height=1.2cm, rounded corners=4pt, align=center, font=\small},
    arr/.style={-{Stealth[length=3mm]}, thick, color=blue!80!black}
]
    % Tầng 1
    \node[layer] (layer1) at (0, 0) {\textbf{TẦNG 1: THIẾT BỊ NHÚNG VÀ CẢM BIẾN IOT (EDGE HARDWARE)}\\ESP32 NodeMCU-32S, DHT22 (Nhiệt/Ẩm), MQ-2 (Gas/Khói), RFID-RC522 (13.56MHz), Relay 4 kênh, Buzzer 5V, Nút nhấn};
    
    % Tầng 2
    \node[layer] (layer2) at (0, 2.0) {\textbf{TẦNG 2: MẠNG VÀ TRUYỀN THÔNG DỮ LIỆU (NETWORK / COMMUNICATION)}\\Giao thức MQTT (Mosquitto Broker QoS 1, Port 1883) \& WebSockets (Socket.IO Real-time Engine)};
    
    % Tầng 3
    \node[layer] (layer3) at (0, 4.0) {\textbf{TẦNG 3: XỬ LÝ NGHIỆP VỤ BACKEND VÀ TRÍ TUỆ NHÂN TẠO (BACKEND SERVICES)}\\NestJS TypeScript Clean Architecture, Command Pattern with Undo, PCCC Rule Engine, Isolation Forest Anomaly Detection, AI Copilot};
    
    % Tầng 4
    \node[layer] (layer4) at (0, 6.0) {\textbf{TẦNG 4: LƯU TRỮ VÀ CƠ SỞ DỮ LIỆU (PERSISTENCE LAYER)}\\Prisma ORM, SQLite / PostgreSQL chuẩn hóa 12 thực thể quan hệ, Lịch sử lệnh, Nhật ký ra vào, Lịch sử PCCC};
    
    % Tầng 5
    \node[layer] (layer5) at (0, 8.0) {\textbf{TẦNG 5: ỨNG DỤNG VÀ TRỰC QUAN HÓA (PRESENTATION / APPLICATION)}\\Next.js 16 App Router: Admin Portal (10 phân hệ giám sát) \& Student Portal (5 phân hệ cá nhân), Xuất báo cáo Excel/CSV};

    % Mũi tên kết nối 2 chiều
    \draw[arr, <->] (layer1) -- (layer2);
    \draw[arr, <->] (layer2) -- (layer3);
    \draw[arr, <->] (layer3) -- (layer4);
    \draw[arr, <->] (layer3) -- (layer5);
\end{tikzpicture}
\caption{Sơ đồ kiến trúc 5 tầng của Hệ thống Quản lý Ký túc xá Thông minh (SmartDorm)}
\label{fig:kien_truc_5_tang}
\end{figure}

Các tầng trong kiến trúc hoạt động phối hợp nhịp nhàng:
\begin{enumerate}
    \item \textbf{Tầng 1 (Edge Hardware):} Thu thập dữ liệu cảm biến vật lý tại từng phòng KTX và điều khiển trực tiếp các tải điện thông qua module relay cách ly quang.
    \item \textbf{Tầng 2 (Communication):} Truyền dẫn dữ liệu nhẹ, tin cậy qua MQTT Broker theo chuẩn QoS 1 và đẩy dữ liệu tức thì lên máy khách qua WebSockets.
    \item \textbf{Tầng 3 (Backend Services):} Trung tâm điều phối toàn bộ luồng nghiệp vụ hệ thống, thực thi xác thực JWT, phân quyền RBAC, kiểm tra điều kiện an toàn PCCC, quản lý hàng đợi hoàn tác lệnh (Undo), và chạy các mô hình AI phát hiện bất thường môi trường.
    \item \textbf{Tầng 4 (Persistence):} Quản lý dữ liệu quan hệ qua Prisma ORM, lưu trữ có cấu trúc nhật ký ra vào, trạng thái thiết bị và các sự cố khẩn cấp.
    \item \textbf{Tầng 5 (Presentation):} Cung cấp trải nghiệm người dùng hiện đại, trực quan, hỗ trợ tương tác trên cả máy tính và thiết bị di động cho ban quản lý và sinh viên.
\end{enumerate}

\section{Công cụ và môi trường triển khai}
\label{sec:cong_cu_moi_truong}

Toàn bộ quá trình nghiên cứu, thiết kế, hiện thực và kiểm nghiệm hệ thống SmartDorm được xây dựng trên nền tảng các công cụ và công nghệ mã nguồn mở tiên tiến nhất. Bảng \ref{tab:cong_cu_trien_khai} tổng hợp danh mục công cụ phần cứng và phần mềm được sử dụng trong dự án.

\begin{table}[H]
\centering
\small
\begin{tabularx}{\textwidth}{|c|p{3.2cm}|p{2.4cm}|X|}
\hline
\textbf{STT} & \textbf{Tên công cụ / Phần mềm} & \textbf{Phiên bản} & \textbf{Vai trò và chức năng sử dụng} \\
\hline
1 & ESP32 NodeMCU-32S & CP2102 Type-C & Vi điều khiển trung tâm nút mạng IoT tại phòng KTX \\
\hline
2 & Cảm biến DHT22 & 3 chân & Đo nhiệt độ (-40 đến 80$^\circ\text{C}$) và độ ẩm (0 đến 100\%) \\
\hline
3 & Cảm biến MQ-2 & 4 chân & Phát hiện nồng độ khí gas dễ cháy (LPG, Propane) và khói \\
\hline
4 & Module Relay 4 kênh & 5V cách ly quang & Đóng/ngắt an toàn phụ tải chiếu sáng, quạt và khóa cửa \\
\hline
5 & Còi Active Buzzer & 5V DC & Phát âm thanh cảnh báo tại chỗ khi có sự cố cháy hoặc thẻ lạ \\
\hline
6 & Nút nhấn 4 chân & Tactile Button & Điều khiển bật/tắt thiết bị tại chỗ và mô phỏng quẹt thẻ RFID \\
\hline
7 & Wokwi Embedded Simulator & Online 2026 & Môi trường mô phỏng sơ đồ mạch và kiểm thử firmware nhúng \\
\hline
8 & PlatformIO / VS Code & Core 6.x & Môi trường biên dịch và nạp firmware C++ cho vi điều khiển \\
\hline
9 & NestJS Framework & v10.x & Nền tảng xây dựng máy chủ backend TypeScript Clean Architecture \\
\hline
10 & Next.js App Router & v16.x & Nền tảng xây dựng giao diện ứng dụng web Admin và Student \\
\hline
11 & Prisma ORM & v5.x & Ánh xạ đối tượng và thao tác cơ sở dữ liệu quan hệ \\
\hline
12 & Eclipse Mosquitto & v2.0.x & MQTT Message Broker trung gian điều phối bản tin IoT \\
\hline
13 & Socket.IO & v4.x & Động cơ WebSockets phát sóng dữ liệu vi khí hậu thời gian thực \\
\hline
14 & Scikit-Learn / Python & v1.4.x & Huấn luyện và vận hành mô hình Isolation Forest phát hiện bất thường \\
\hline
15 & Telegram Bot API & BotFather 2026 & Kênh gửi thông báo khẩn cấp PCCC có nút tương tác xác nhận \\
\hline
16 & Nodemailer SMTP & v6.9.x & Dịch vụ gửi thư điện tử cảnh báo khẩn cấp định dạng HTML \\
\hline
\end{tabularx}
\caption{Danh sách công cụ, phần mềm và môi trường triển khai hệ thống}
\label{tab:cong_cu_trien_khai}
\end{table}

\newpage
